\section{Nonflat norms without potential theory on curves}\label{sec:nf-route}

This chapter proves the nonflat norm statement of Proposition~\ref{prop:nonflat-norm} for the functions to which the
local regularization theorem is applied in Corollary~\ref{cor:normal-weak-global}, namely functions whose pullback to
the charts of a blowup is regularizable (Definition~\ref{def:nf-chart-regularizable}). The weights are defined by
characteristic polynomials at a rigid lift (Definition~\ref{def:nf-weights}). The identification of the extended norm
off the \'etale locus (the last step of the proof of \cite[Proposition~5.22]{FHX}) is made on a line through the rigid
lift, using Gauss points of small discs, the hull inequality on $S$ and a weighted maximum principle for pullbacks of
Fubini--Study functions along a finite free cover of the line (Lemma~\ref{lem:nf-pieces}). Subharmonic functions in the
sense of \cite{GR25} and Lemma~\ref{lem:curves} are not used; see Remark~\ref{rmk:nf-source}.

\subsection{Setting and the function class}\label{subsec:nf-setting}

\begin{convention}\label{conv:nf-setting}
\uses{conv:sec5, prop:nonflat-norm}
As in Convention~\ref{conv:sec5}, $k$ is countable, algebraically closed and trivially valued. Throughout,
$Y=\Spec A$ is a normal affine variety of dimension $n\ge1$, $S=\Spec B$ with $B\cong k[x_1,\dots,x_n]_{h_S}$
($0\ne h_S$) a principal open subset of $\A^n$, $\varphi\colon B\to A$ is finite and injective, $\pi=\Spec\varphi$ is
generically \'etale of degree $d=[\mathrm{Frac}\,A:\mathrm{Frac}\,B]$, and $Z=V(I_Z)\subsetneq Y$ is closed, $I_Z\ne0$.
For a finite morphism $\Spec R'\to\Spec Q$ and a point $z$ of $(\Spec Q)^\an$, the points over $z$ correspond to the
maximal ideals of the finite $\mathcal H(z)$-algebra $\mathcal H(z)\otimes_QR'$, and the multiplicity $m_x$ of a point
$x$ over $z$ is the dimension over $\mathcal H(z)$ of the local factor of $\mathcal H(z)\otimes_QR'$ at $x$
(Lemma~\ref{lem:finite-flat-norm}). For a function $w$ on $(\Spec R')^\an$ we write
$\Norm w(z)=\sum_{x\mapsto z}m_xw(x)$. If $F$ is a complete valued field extension of $k$ and $R$ is an $F$-algebra,
the points of $(\Spec R)^\an$ over $F$ are the multiplicative seminorms on $R$ restricting to the absolute value of $F$.
\end{convention}

\begin{definition}[Functions regularizable on the charts of a blowup]\label{def:nf-chart-regularizable}
\lean{EnvelopeProperty.IsChartRegularizable, EnvelopeProperty.IsRegularizable.isChartRegularizable}
\leanok
\uses{def:fs, def:weakly-psh, rmk:weakly-psh-blowup}
A function $u\colon Y^\an\to[-\infty,\infty)$ is \emph{regularizable on the charts of a blowup} if there is a finite set
$s\subseteq A$ containing a nonzero element such that for every $\sigma\in s\setminus\{0\}$ the function
$u\circ\mu_\sigma$ lies in $\Rcal(A[s/\sigma])$, where $A[s/\sigma]\subseteq\mathrm{Frac}\,A$ is the chart of the blowup
of $Y$ along $(s)$ and $\mu_\sigma\colon(\Spec A[s/\sigma])^\an\to Y^\an$ is restriction along $A\subseteq A[s/\sigma]$.
The set $s$ is called a \emph{witness}.
\end{definition}

\begin{lemma}[Closure properties of the class]\label{lem:nf-chart-regularizable-closure}
\lean{EnvelopeProperty.isPshClass_isChartRegularizable, EnvelopeProperty.IsChartRegularizable.isWeaklyPSHOn, EnvelopeProperty.IsChartRegularizable.upperSemicontinuous, EnvelopeProperty.IsChartRegularizable.add_isRegularizable, EnvelopeProperty.IsChartRegularizable.max_isRegularizable, EnvelopeProperty.IsChartRegularizable.add_const, EnvelopeProperty.IsChartRegularizable.comp_awayComap, EnvelopeProperty.IsChartRegularizable.comp_proj_add_logAbs_X}
\leanok
\uses{def:nf-chart-regularizable, lem:reg-closure-elementary, lem:reg-triv-finite, lem:fs-basic, def:weakly-psh, rmk:weakly-psh-blowup}
Let $u$ be regularizable on the charts of a blowup, with witness $s$.
\begin{enumerate}
\item $u$ is weakly psh with the blowup along $(s)$ as witness; in particular $u$ is upper semicontinuous.
\item For $v\in\Rcal(Y)$, the functions $u+v$ and $\max\{u,v\}$ are regularizable on the charts of a blowup, with
witness $s$.
\item For $0\ne h\in A$, the restriction of $u$ to $Y_h^\an$ is regularizable on the charts of a blowup of $Y_h$,
with witness the image of $s$.
\item On $\widetilde Y=Y\times\A^1=\Spec A[w]$ with projection $p$, the function $u\circ p+\log|w|$ is regularizable on
the charts of a blowup, with witness $s\subseteq A\subseteq A[w]$.
\end{enumerate}
\end{lemma}

\begin{proof}
(1) Every regularizable function is a decreasing limit of Fubini--Study functions, which are convex PL
(Lemma~\ref{lem:fs-basic}); so $u\circ\mu_\sigma$ is locally psh on every chart, which is
Definition~\ref{def:weakly-psh} in the form of Remark~\ref{rmk:weakly-psh-blowup}; weakly psh functions are upper
semicontinuous. (2) The function $v\circ\mu_\sigma$ is finite at the trivial valuation of $A[s/\sigma]$, which lies over
the trivial valuation of $A$ (Lemma~\ref{lem:reg-triv-finite}), so $v\circ\mu_\sigma\in\Rcal(A[s/\sigma])$ by
Lemma~\ref{lem:reg-closure-elementary}(4). Sums and maxima of regularizable functions are regularizable
(Lemma~\ref{lem:reg-closure-elementary}), and $(u+v)\circ\mu_\sigma=u\circ\mu_\sigma+v\circ\mu_\sigma$,
$\max\{u,v\}\circ\mu_\sigma=\max\{u\circ\mu_\sigma,v\circ\mu_\sigma\}$. (3) Inside $\mathrm{Frac}\,A$, the chart
$A_h[s/\sigma]$ of the blowup of $Y_h$ is generated by $A[s/\sigma]$ and $h^{-1}$; the restriction of $u|_{Y_h^\an}$ to
its chart is $u\circ\mu_\sigma$ pulled back along $A[s/\sigma]\to A_h[s/\sigma]$, which is regularizable by
Lemma~\ref{lem:reg-closure-elementary}(4) (it is finite at the trivial valuation, which lies over the trivial valuation
of $A[s/\sigma]$). (4) The chart of the blowup of $\widetilde Y$ along $(s)$ is $A[w][s/\sigma]=A[s/\sigma][w]$. On it,
$u\circ p$ is the pullback of $u\circ\mu_\sigma$ along $A[s/\sigma]\to A[s/\sigma][w]$, regularizable by
Lemma~\ref{lem:reg-closure-elementary}(4), and $\log|w|$ is regularizable by Lemma~\ref{lem:reg-closure-elementary}(3);
their sum is regularizable.
\end{proof}

\begin{lemma}[The functions of the descent corollary]\label{lem:nf-descent-chart}
\lean{EnvelopeProperty.ProjSubvariety.BlowupCenter.isChartRegularizable_descentChartFun}
\leanok
\uses{def:nf-chart-regularizable, lem:pg-descent-weakly-psh}
In the setting of Theorem~\ref{thm:pg-psh-descent}, for every $i$ with $x_i\notin I$ the chart function $u_i$ on
$Y_i^\an$ is regularizable on the charts of a blowup, with witness $\mathfrak a_i$.
\end{lemma}

\begin{proof}
The proof of Lemma~\ref{lem:pg-descent-weakly-psh} shows that on every chart
$\Spec A_{x_i}[\mathfrak a_i/g_{il}]$, $g_{il}\ne0$, the function $w\mapsto u_i(w|_{A_{x_i}})$ is the limit of a
decreasing sequence of Fubini--Study functions of the chart, and it is finite at the trivial valuation; this is
Definition~\ref{def:nf-chart-regularizable}.
\end{proof}

\subsection{Characteristic polynomials over a normal base}\label{subsec:nf-charpoly}

\begin{lemma}[Characteristic polynomials of integral elements]\label{lem:nf-charpoly}
\lean{EnvelopeProperty.nonflatCharpoly, EnvelopeProperty.charpoly_lmul_tensor_eq_map_nonflatCharpoly, EnvelopeProperty.natDegree_nonflatCharpoly, EnvelopeProperty.aeval_nonflatCharpoly, EnvelopeProperty.dvd_algebraMap_coeff_zero_nonflatCharpoly, EnvelopeProperty.coeff_zero_nonflatCharpoly_ne_zero, EnvelopeProperty.nonflatCharpoly_comparison, EnvelopeProperty.nonflat_exists_free_localizationAway}
\leanok
Let $R$ be an integrally closed domain with fraction field $K$, and $T$ a domain containing $R$ and finite over $R$,
with $\dim_K(K\otimes_RT)=d$. For $b\in T$ let $\chi_b\in K[X]$ be the characteristic polynomial of multiplication by
$b$ on the $K$-vector space $K\otimes_RT$.
\begin{enumerate}
\item $\chi_b\in R[X]$, and $\chi_b$ is monic of degree $d$.
\item $\chi_b(b)=0$ in $T$.
\item $\chi_b(0)\in bT$, and $\chi_b(0)\ne0$ if $b\ne0$.
\item Let $0\ne f\in R$ be such that $T_f$ is a free $R_f$-module, let $\psi\colon R\to Q$ be a homomorphism to a domain
with $\psi(f)\ne0$, let $T'$ be a finite free $Q$-algebra and $\theta\colon T\to T'$ a homomorphism compatible with
$\psi$, and assume that for some $0\ne c\in Q$ the induced map $Q_{c\psi(f)}\otimes_{R_f}T_f\to T'_{c\psi(f)}$ is an
isomorphism. Then the characteristic polynomial of multiplication by $\theta(b)$ on $T'$ over $Q$ is $\psi(\chi_b)$.
\end{enumerate}
\end{lemma}

\begin{proof}
(1) $L=K\otimes_RT=\mathrm{Frac}\,T$ is a field, finite over $K$, and $\chi_b=\mu_b^{[L:K(b)]}$ for the minimal
polynomial $\mu_b$ of $b$ over $K$. Since $b$ is integral over the integrally closed domain $R$, $\mu_b\in R[X]$.
(2) By the Cayley--Hamilton theorem $\chi_b(b)=0$ in $L$, and $T\subseteq L$. (3) By (2),
$\chi_b(0)=-b\,(b^{d-1}+c_{d-1}b^{d-2}+\dots+c_1)$ with $c_i\in R$; and $\pm\chi_b(0)$ is the norm of $b$ from $L$ to
$K$, nonzero when $b\ne0$. (4) Since $T_f$ is free over $R_f$ and $K\otimes_{R_f}T_f=L$, the characteristic polynomial
of $b$ on $T_f$ over $R_f$ is $\chi_b$; characteristic polynomials of free modules commute with base change, so the
characteristic polynomial of $\theta(b)$ on $Q_{c\psi(f)}\otimes_{R_f}T_f\cong T'_{c\psi(f)}$ over $Q_{c\psi(f)}$ is
$\psi(\chi_b)$. It is also the image of the characteristic polynomial $\chi'$ of $\theta(b)$ on $T'$ over $Q$, by base
change along $Q\to Q_{c\psi(f)}$. This map is injective because $Q$ is a domain and $c\psi(f)\ne0$, so
$\chi'=\psi(\chi_b)$.
\end{proof}

\begin{lemma}[Characteristic polynomials on local Artinian factors]\label{lem:nf-local-charpoly}
\lean{EnvelopeProperty.charpoly_lmul_eq_prod_nonflatResidue, EnvelopeProperty.nonflat_rootMultiplicity_charpoly_lmul, EnvelopeProperty.nonflatResidue, EnvelopeProperty.nonflat_finrank_localization_eq_length}
\leanok
\uses{spectrum:algclosed-local-norm-form}
Let $F$ be algebraically closed and $E$ a finite $F$-algebra, $E=\prod_xE_x$ its decomposition into local factors
with residue fields $F$, and $e_x=\dim_FE_x$. For $b\in E$ with residues $b(x)\in F$, the characteristic polynomial of
multiplication by $b$ on $E$ is $\prod_x(X-b(x))^{e_x}$.
\end{lemma}

\begin{proof}
Multiplication by $b$ preserves each factor $E_x$, so the characteristic polynomial is the product of those on the
$E_x$. On $E_x$, $b-b(x)$ lies in the maximal ideal of the Artinian local ring $E_x$, hence is nilpotent, so its
characteristic polynomial is $X^{e_x}$ and that of $b$ is $(X-b(x))^{e_x}$.
\end{proof}

\subsection{The good locus, rigid data and the weights}\label{subsec:nf-weights}

\begin{lemma}[A good principal open of $S$]\label{lem:nf-good-locus}
\lean{EnvelopeProperty.IsNonflatGoodLocus, EnvelopeProperty.exists_isNonflatGoodLocus, EnvelopeProperty.IsNonflatGoodLocus.exists_mem_ne_zero, EnvelopeProperty.IsNonflatGoodLocus.awayMap_etale, EnvelopeProperty.IsNonflatGoodLocus.awayMap_free, EnvelopeProperty.nonflat_awayMap_finite}
\leanok
\uses{conv:nf-setting, lem:nf-charpoly, lem:finite-flat-norm}
There is $0\ne h\in B$ such that $A_{\varphi(h)}$ is a finite free \'etale $B_h$-algebra and
$Y_0=\Spec A_{\varphi(h)}=\pi^{-1}(S_h)$ is contained in $Y\setminus Z$.
\end{lemma}

\begin{proof}
By generic \'etaleness there is $0\ne h_1\in B$ with $B_{h_1}\to A_{\varphi(h_1)}$ \'etale. The finite module $A$ over
the Noetherian domain $B$ is free at the generic point, so there is $0\ne h_2\in B$ with $A_{\varphi(h_2)}$ free over
$B_{h_2}$ (as in the proof of Lemma~\ref{lem:finite-flat-norm}). Choose $0\ne f\in I_Z$ and put $c_f=\chi_f(0)$
(Lemma~\ref{lem:nf-charpoly} for $B\subseteq A$); then $0\ne c_f\in B$ and $\varphi(c_f)\in fA\subseteq I_Z$. Put
$h=h_1h_2c_f$. Localization at a further element preserves finiteness, freeness and \'etaleness, and $\varphi(h)$
vanishes on $Z$, so $Y_0\cap Z=\varnothing$.
\end{proof}

We fix $h$ as in Lemma~\ref{lem:nf-good-locus}, $S_0=S_h$ and $Y_0=Y_{\varphi(h)}$.

\begin{definition}[Rigid data]\label{def:nf-rigid-data}
\lean{EnvelopeProperty.NonflatRigidEmbedding, EnvelopeProperty.NonflatRigidEmbedding.basePoint, EnvelopeProperty.NonflatRigidEmbedding.fibre, EnvelopeProperty.nonflatRigidField, EnvelopeProperty.nonflatRigidEmbedding, EnvelopeProperty.nonflatRigidField_norm_algebraMap, EnvelopeProperty.nonflatRigidField_exists_one_lt_norm}
\leanok
\uses{spectrum:gauss-point, spectrum:completed-residue-field-pullback, spectrum:complete-valued-algebraic-closure}
Let $s\in S^\an$. A \emph{rigid datum} for $s$ is a complete, algebraically closed, nontrivially valued field $F$ with
an isometric embedding $\mathcal H(s)\to F$. The composite $B\to\mathcal H(s)\to F$ extends to the $F$-point
$\xi\colon B_F=F\otimes_kB\to F$, and we write $A_F=F\otimes_kA$, $\pi_F\colon\Spec A_F\to\Spec B_F$, and
$\pi_F^{-1}(\xi)$ for the set of points of $(\Spec A_F)^\an$ over $F$ lying over $\xi$. A rigid datum exists: let $s'$
be the Gauss point of radius $1/2$ over $s$ on $\Spec B[w]$; then $\mathcal H(s)\to\mathcal H(s')$ is isometric and
$|w|_{s'}=1/2$, and the completed algebraic closure of $\mathcal H(s')$ is complete, algebraically closed and
nontrivially valued.
\end{definition}

\begin{lemma}[Rigid fibres]\label{lem:nf-rigid-fibre}
\lean{EnvelopeProperty.nonflat_baseChange_finite, EnvelopeProperty.NonflatRigidEmbedding.finite_fibre, EnvelopeProperty.NonflatRigidEmbedding.fibreEval, EnvelopeProperty.NonflatRigidEmbedding.val_eq_nnnorm_fibreEval, EnvelopeProperty.NonflatRigidEmbedding.exists_mem_fibre_comap_includeRight_eq, EnvelopeProperty.NonflatRigidEmbedding.surjOn_comap_includeRight_fibre, EnvelopeProperty.NonflatRigidEmbedding.exists_fibreEval_injective, EnvelopeProperty.nonflat_exists_injOn_algHom, EnvelopeProperty.nonflatBaseEquiv, EnvelopeProperty.nonflat_isDomain_baseChange_base, EnvelopeProperty.nonflat_isIntegrallyClosed_baseChange_base, EnvelopeProperty.IsNonflatGoodLocus.finrank_baseChange}
\leanok
\uses{def:nf-rigid-data, spectrum:finite-completed-fibres, spectrum:complete-norm-extension, spectrum:finite-rational-rigid-fibre-lift, algebra:geometric-domain-base-change, algebra:principal-open-scalar-tensor}
Let $(F,\xi)$ be a rigid datum for $s\in S^\an$.
\begin{enumerate}
\item $B_F\cong F[x_1,\dots,x_n]_{h_S}$ is an integrally closed domain, $A_F$ is a domain, finite over $B_F$, and
$\dim_{\mathrm{Frac}\,B_F}(\mathrm{Frac}\,B_F\otimes_{B_F}A_F)=d$.
\item $\pi_F^{-1}(\xi)$ is finite. Each $z\in\pi_F^{-1}(\xi)$ is an $F$-point: there is a homomorphism
$A_F\to F$, $b\mapsto b(z)$, over $\xi$, with $|b|_z=|b(z)|$.
\item The restriction map $r\colon\pi_F^{-1}(\xi)\to\pi^{-1}(s)$ is surjective.
\item There is $b_s\in A_F$ whose values $b_s(z)$, $z\in\pi_F^{-1}(\xi)$, are pairwise distinct.
\end{enumerate}
\end{lemma}

\begin{proof}
(1) $B_F$ is a localization of a polynomial ring over a field; $A_F$ is a domain because $k$ is algebraically closed
(Lemma~\ref{algebra:geometric-domain-base-change}); finiteness and the degree follow by base change from a free
localization $A_{\varphi(h)}$ of rank $d$ over $B_h$ (Lemma~\ref{lem:nf-good-locus}), as $\mathrm{Frac}\,B_F$ is a
$B_h$-algebra. (2) The points over $\xi$ correspond to the maximal ideals of the finite $F$-algebra
$F\otimes_{B_F,\xi}A_F$; their residue fields are finite over $F$, hence equal to $F$, and the absolute value on each is
the unique extension of that of $F$. (3) Let $y\in\pi^{-1}(s)$. The completed residue field $\mathcal H(y)$ is finite
over $\mathcal H(s)$, so there is an $\mathcal H(s)$-embedding $\tau\colon\mathcal H(y)\to F$; it is isometric, since the
absolute value of the complete field $\mathcal H(s)$ has a unique extension to $\mathcal H(y)$. The homomorphism
$F\otimes_kA\to F$, $c\otimes a\mapsto c\,\tau(a)$, lies over $\xi$ and defines $z\in\pi_F^{-1}(\xi)$ with $r(z)=y$.
(4) The map $A_F\to\prod_zF$, $b\mapsto(b(z))_z$, is surjective by the Chinese remainder theorem, since the $z$ have
pairwise distinct kernels; take a preimage of a vector with distinct entries.
\end{proof}

\begin{definition}[The weights]\label{def:nf-weights}
\lean{EnvelopeProperty.nonflatBaseCharpoly, EnvelopeProperty.NonflatRigidEmbedding.rigidMult, EnvelopeProperty.NonflatRigidEmbedding.rigidWeight, EnvelopeProperty.nonflatSeparating, EnvelopeProperty.nonflatWeight, EnvelopeProperty.nonflatNorm}
\leanok
\uses{lem:nf-good-locus, def:nf-rigid-data, lem:nf-rigid-fibre, lem:nf-charpoly, lem:finite-etale-norm}
For $s\in S_0^\an$ and $y\in\pi^{-1}(s)$, let $m_y(\pi)$ be the multiplicity of Lemma~\ref{lem:finite-etale-norm}.
For $s\in S^\an\setminus S_0^\an$, fix a rigid datum $(F_s,\xi_s)$ and $b_s$ as in Lemma~\ref{lem:nf-rigid-fibre}(4),
and let $\chi_s\in B_{F_s}[X]$ be the characteristic polynomial of $b_s$ over $\mathrm{Frac}\,B_{F_s}$
(Lemma~\ref{lem:nf-charpoly}). For $z\in\pi_{F_s}^{-1}(\xi_s)$ let $M_z$ be the multiplicity of $b_s(z)$ as a root of
$\chi_s(\xi_s)\in F_s[X]$, the polynomial obtained by applying $\xi_s$ to the coefficients. For $y\in\pi^{-1}(s)$ put
$m_y(\pi)=\sum_{z\in\pi_{F_s}^{-1}(\xi_s),\,r(z)=y}M_z$. For a function $u$ on $Y^\an$ put
$\Norm_\pi u(s)=\sum_{\pi(y)=s}m_y(\pi)u(y)$.
\end{definition}

The weights depend on the choices made for each $s$, but not on any function $u$.

\subsection{Lines through a rigid lift}\label{subsec:nf-lines}

\begin{definition}[Admissible lines and their normalized curves]\label{def:nf-line}
\uses{def:nf-rigid-data, algebra:affine-line-substitution, algebra:principal-open-line-ring, algebra:line-preimage-actual-components, algebra:normalized-line-component-ring}
Let $(F,\xi)$ be a rigid datum for $s\in S^\an$, with $\xi$ the point $a\in F^n$, $h_S(a)\ne0$. For a direction
$0\ne w\in F^n$ let $\iota\colon B_F\to Q$ be the restriction to the line $t\mapsto a+tw$, where
$Q=F[t]_{p}$, $p(t)=h_S(a+tw)$; then $p(0)\ne0$, $Q$ is a principal ideal domain, and composing $\iota$ with evaluation
at $t=0$ gives $\xi$. Put $A_C=Q\otimes_{B_F}A_F$ and let $R'=\prod_PD_P$, the product over the minimal primes $P$ of
$A_C$ of the integral closures $D_P$ of $A_C/P$ in their fraction fields; let $q\colon\Spec R'\to\Spec Q$ and
$p\colon\Spec R'\to\Spec A_F$ be the natural maps. The direction $w$ is \emph{admissible} for an element
$0\ne e\in B$ if $\iota(e)\ne0$ in $Q$.
\end{definition}

\begin{lemma}[Admissible directions exist]\label{lem:nf-line-exists}
\lean{EnvelopeProperty.nonflatLinePoint, EnvelopeProperty.nonflat_lineEval_comp, EnvelopeProperty.nonflat_exists_admissible_direction, EnvelopeProperty.nonflat_admissible_of_mul}
\leanok
\uses{def:nf-line, algebra:affine-line-avoids-polynomial}
For every $0\ne e\in B$ and every rigid datum there is a direction admissible for $e$.
\end{lemma}

\begin{proof}
Write $e=e^\circ/h_S^m$ with $0\ne e^\circ\in k[x]$. Then $\iota(e)\ne0$ if and only if $e^\circ(a+tw)\ne0$ in $F[t]$,
and since $F$ is infinite such a $w$ exists (Lemma~\ref{algebra:affine-line-avoids-polynomial}).
\end{proof}

\begin{lemma}[The normalized curve of an admissible line]\label{lem:nf-line-curve}
\lean{EnvelopeProperty.AffineLine.nonflatLineTensorEquiv, EnvelopeProperty.AffineLine.nonflatLine_exists_baseChange_bijective, EnvelopeProperty.AffineLine.nonflatLine_exists_mul_mem_range, EnvelopeProperty.AffineLine.nonflatLine_exists_pow_mul_eq_zero, EnvelopeProperty.AffineLine.nonflatLine_finrank_eq, EnvelopeProperty.AffineLine.nonflatLine_finrank_eq_finrank, EnvelopeProperty.AffineLine.nonflatLine_normalizedCurve_structure, EnvelopeProperty.AffineLine.nonflatLine_tensor_etale}
\leanok
\uses{def:nf-line, lem:nf-good-locus, lem:nf-rigid-fibre, algebra:finite-free-normalized-line-product, algebra:line-component-domination, algebra:line-tensor-original-unreduced-preimage, algebra:line-tensor-free-rank-etale}
Let $w$ be admissible for $h$. Then:
\begin{enumerate}
\item each $D_P$ is a Dedekind domain, finite over $Q$, with $Q\to D_P$ injective, and $R'$ is a finite free $Q$-module;
\item there is $0\ne c\in Q$ such that $A_C\to R'$ becomes an isomorphism after inverting $c$;
\item $R'$ has rank $d$ over $Q$.
\end{enumerate}
\end{lemma}

\begin{proof}
(1) This is Lemma~\ref{algebra:finite-free-normalized-line-product} and the domination of the components of $A_C$
(Lemma~\ref{algebra:line-component-domination}). (2) Let $K=\mathrm{Frac}\,Q$. Since $\iota(h)\ne0$, $K$ is a
$B_{F,h}$-algebra and $K\otimes_QA_C=K\otimes_{B_{F,h}}A_{F,h}$ is a finite \'etale $K$-algebra of dimension $d$
(Lemma~\ref{algebra:line-tensor-free-rank-etale}), hence a finite product of fields. Every minimal prime $P$ of $A_C$
lies over $(0)$ of $Q$ (Lemma~\ref{algebra:line-component-domination}), so the minimal primes of $K\otimes_QA_C$ are the
$K\otimes P$, and $K\otimes_QA_C\cong\prod_PK\otimes_Q(A_C/P)$, the nilradical of $A_C$ vanishing after
$\otimes K$. Each $K\otimes_Q(A_C/P)$ is a domain finite over $K$, hence equal to $\mathrm{Frac}(A_C/P)$; and
$K\otimes_QD_P=\mathrm{Frac}(D_P)=\mathrm{Frac}(A_C/P)$ as $D_P$ is finite over $Q$. So $K\otimes A_C\to K\otimes R'$ is
an isomorphism. The kernel and the cokernel of $A_C\to R'$ are finitely generated $Q$-modules ($A_C$ and $R'$ are finite
over the Noetherian ring $Q$) that vanish after $\otimes K$; they are annihilated by a single $0\ne c\in Q$, and
localization at $c$ is exact. (3) $\dim_K(K\otimes R')=\dim_K(K\otimes A_C)=d$.
\end{proof}

\begin{lemma}[Line multiplicities are the weights]\label{lem:nf-line-weights}
\lean{EnvelopeProperty.AffineLine.nonflatLine_lineWeights, EnvelopeProperty.AffineLine.nonflatLine_charpoly_lmul_eq, EnvelopeProperty.nonflat_charpoly_map_eq_prod_fibre, EnvelopeProperty.nonflat_sum_flatFibreMult_eq_rootMultiplicity, EnvelopeProperty.nonflat_sum_rootMultiplicity_eq_finrank, EnvelopeProperty.nonflat_one_le_rootMultiplicity_map_nonflatCharpoly, EnvelopeProperty.nonflat_one_le_rootMultiplicity_of_mem_lineBaseFibre, EnvelopeProperty.nonflat_exists_comap_eq_of_mem_lineBaseFibre, EnvelopeProperty.nonflat_flatNormFun_comap_eq_sum_rootMultiplicity}
\leanok
\uses{def:nf-weights, lem:nf-line-curve, lem:nf-charpoly, lem:nf-local-charpoly, lem:nf-rigid-fibre}
Let $s\notin S_0^\an$, with the data of Definition~\ref{def:nf-weights}, and let $w$ be admissible for $h$. Let $\xi$
also denote the point $t=0$ of $(\Spec Q)^\an$. Then every point $x$ of $(\Spec R')^\an$ over $\xi$ is an $F$-point,
$p(x)\in\pi_F^{-1}(\xi)$, and for every $z\in\pi_F^{-1}(\xi)$
\[
\sum_{x\mapsto\xi,\ p(x)=z}m_x(q)=M_z .
\]
Moreover $M_z\ge1$ and $\sum_zM_z=d$.
\end{lemma}

\begin{proof}
The points over $\xi$ correspond to the maximal ideals of the finite $F$-algebra $F\otimes_{Q,\xi}R'$, whose residue
fields are $F$; the restriction of $x$ to $A_F$ lies over $\iota^*\xi=\xi$. Let $\chi'$ be the characteristic
polynomial of (the image of) $b_s$ on the free $Q$-module $R'$. By Lemma~\ref{lem:nf-charpoly}(4), with $R=B_F$,
$T=A_F$, $f=h$, $\psi=\iota$, $T'=R'$ and $c$ from Lemma~\ref{lem:nf-line-curve}(2) (note $A_{F,h}$ is free over
$B_{F,h}$ and $Q_{c\iota(h)}\otimes_{B_{F,h}}A_{F,h}=(A_C)_{c\iota(h)}$), $\chi'=\iota(\chi_s)$. Applying evaluation at
$t=0$ to the coefficients, $\iota(\chi_s)$ becomes $\chi_s(\xi)$. On the other hand, by base change of the free module
$R'$, the image of $\chi'$ is the characteristic polynomial of $b_s$ on $F\otimes_{Q,\xi}R'$, which is
$\prod_x(X-b_s(x))^{m_x(q)}$ by Lemma~\ref{lem:nf-local-charpoly}, and $b_s(x)=b_s(p(x))$. Hence
$\prod_x(X-b_s(p(x)))^{m_x(q)}=\chi_s(\xi)$. As the $b_s(z)$ are pairwise distinct, comparing the multiplicity of the
root $b_s(z)$ gives the displayed formula. By Lemma~\ref{lem:nf-charpoly}(2), $\chi_s(b_s)=0$ in $A_F$; applying the
point $z$ gives $\chi_s(\xi)(b_s(z))=0$, so $M_z\ge1$. Every root of $\chi_s(\xi)$ is some $b_s(p(x))$, so
$\sum_zM_z=\deg\chi_s=d$.
\end{proof}

\begin{lemma}[Basic properties of the weights]\label{lem:nf-weights-basic}
\lean{EnvelopeProperty.one_le_nonflatWeight, EnvelopeProperty.sum_nonflatWeight, EnvelopeProperty.NonflatRigidEmbedding.one_le_rigidMult, EnvelopeProperty.nonflatNormWeights}
\leanok
\uses{def:nf-weights, lem:nf-line-weights, lem:nf-line-exists, lem:nf-rigid-fibre, lem:finite-etale-norm, lem:nf-whole-rigid-weighted-norm}
For every $s\in S^\an$ the fibre $\pi^{-1}(s)$ is finite, $m_y(\pi)\ge1$ for $y\in\pi^{-1}(s)$, and
$\sum_{\pi(y)=s}m_y(\pi)=d$.
\end{lemma}

\begin{proof}
Fibres of finite maps are finite. For $s\in S_0^\an$ this is Lemma~\ref{lem:finite-etale-norm}. Otherwise, choose a
direction admissible for $h$ (Lemma~\ref{lem:nf-line-exists}); by Lemma~\ref{lem:nf-line-weights} every $M_z\ge1$ and
$\sum_zM_z=d$. By Lemma~\ref{lem:nf-rigid-fibre}(3) every $y\in\pi^{-1}(s)$ is $r(z)$ for some $z$, so $m_y(\pi)\ge1$,
and $\sum_ym_y(\pi)=\sum_zM_z=d$.
\end{proof}

\subsection{Gauss rays and generic fibres}\label{subsec:nf-gauss}

\begin{lemma}[Gauss rays on a line]\label{lem:nf-gauss-ray}
\lean{EnvelopeProperty.nonflatFieldPoint, EnvelopeProperty.nonflatEvalPoint, EnvelopeProperty.nonflatGaussPoint, EnvelopeProperty.nonflatGaussPoint_zero, EnvelopeProperty.continuous_nonflatGaussPoint, EnvelopeProperty.nonflatGaussPoint_val_ne_zero, EnvelopeProperty.tendsto_nonflatGaussPoint_zero, EnvelopeProperty.exists_nonflatGaussPoint_maxModulus, EnvelopeProperty.eventually_nonflatGaussPoint_maxModulus}
\leanok
\uses{spectrum:centered-gauss-whole-ray, spectrum:centered-gauss-generic-support, spectrum:centered-gauss-principal-open-lift}
Let $F$ be a complete valued field and $Q=F[t]_p$ with $p(0)\ne0$. For $\epsilon\ge0$ let $\zeta_\epsilon$ be the point
of $(\Spec Q)^\an$ with $|\sum_ia_it^i|_{\zeta_\epsilon}=\max_i|a_i|\epsilon^i$ on $F[t]$. Then $\zeta_0$ is evaluation at
$0$, $\epsilon\mapsto\zeta_\epsilon$ is continuous, $\zeta_\epsilon$ has zero kernel for $\epsilon>0$, and there is
$\epsilon_0>0$ such that $|f(0)|\le|f|_{\zeta_\epsilon}$ for all $f\in Q$ and $0<\epsilon<\epsilon_0$.
\end{lemma}

\begin{proof}
The first three assertions are Lemmas~\ref{spectrum:centered-gauss-whole-ray}, \ref{spectrum:centered-gauss-generic-support}
and~\ref{spectrum:centered-gauss-principal-open-lift} (the ray lies in the principal open because $p(0)\ne0$). Choose
$\epsilon_0>0$ with $|p_i|\epsilon_0^i<|p(0)|$ for the finitely many coefficients $p_i$, $i\ge1$; then
$|p|_{\zeta_\epsilon}=|p(0)|$ for $\epsilon<\epsilon_0$. For $f=P/p^m$ with $P\in F[t]$, the constant term gives
$|P(0)|\le|P|_{\zeta_\epsilon}$, so $|f(0)|=|P(0)|/|p(0)|^m\le|P|_{\zeta_\epsilon}/|p|_{\zeta_\epsilon}^m=|f|_{\zeta_\epsilon}$.
\end{proof}

\begin{lemma}[Approach through the \'etale locus]\label{lem:nf-hull-approach}
\lean{EnvelopeProperty.nonflat_comap_nonflatGaussPoint_val_ne_zero, EnvelopeProperty.tendsto_comap_nonflatGaussPoint, EnvelopeProperty.eventually_mem_hull_comap_nonflatGaussPoint, EnvelopeProperty.IsRegularizable.nonflat_hull_approach, EnvelopeProperty.NonflatRigidEmbedding.comap_coordLine_nonflatEvalPoint}
\leanok
\uses{lem:nf-gauss-ray, def:nf-line, lem:hull-basic, def:hull}
Let $s\in S^\an\setminus S_0^\an$, $(F,\xi)$ a rigid datum and $w$ admissible for $h$. Put
$s_\epsilon=(\iota^*\zeta_\epsilon)|_B\in S^\an$. Then $s_\epsilon\in S_0^\an$ for $\epsilon>0$, $s_\epsilon\to s$ as
$\epsilon\to0$, and $s\in\widehat{\{s_\epsilon\}}$ for $0<\epsilon<\epsilon_0$. Consequently, for every $v\in\Rcal(S)$,
$v(s)\le v(s_\epsilon)$ for $0<\epsilon<\epsilon_0$ and $\limsup_{\epsilon\to0}v(s_\epsilon)\le v(s)$.
\end{lemma}

\begin{proof}
$|h|_{s_\epsilon}=|\iota(h)|_{\zeta_\epsilon}\ne0$, since $\iota(h)\ne0$ and $\zeta_\epsilon$ has zero kernel
(Lemma~\ref{lem:nf-gauss-ray}). Restriction maps are continuous and $s_0=\xi|_B=s$, so $s_\epsilon\to s$. For $f\in B$,
$|f|_s=|\iota(f)(0)|\le|\iota(f)|_{\zeta_\epsilon}=|f|_{s_\epsilon}$ by Lemma~\ref{lem:nf-gauss-ray}, i.e.
$s\in\widehat{\{s_\epsilon\}}$. The hull inequality for $v$ (Lemma~\ref{lem:hull-basic}) with the compact set
$\{s_\epsilon\}$ gives $v(s)\le v(s_\epsilon)$, and upper semicontinuity of $v$ gives the last assertion.
\end{proof}

\begin{lemma}[Generic fibres of the normalized curve]\label{lem:nf-generic-fibre}
\lean{EnvelopeProperty.NonflatRigidEmbedding.exists_flatNormFun_line_eq_nonflatNorm, EnvelopeProperty.AffineLine.nonflatLine_exists_flatNormFun_eq, EnvelopeProperty.AffineLine.nonflatLine_exists_baseChangeAlgHom_bijective, EnvelopeProperty.nonflat_flatNormFun_comap_eq_of_bijective, EnvelopeProperty.nonflat_flatNormFun_comap_eq_of_baseChange_bijective, EnvelopeProperty.Berk.nonflat_flatNormFun_comap_of_tensorMap_bijective, EnvelopeProperty.nonflatBaseChangeAlgHom}
\leanok
\uses{lem:nf-line-curve, def:nf-weights, spectrum:finite-base-change-multiplicity-mass, spectrum:all-function-finite-base-change, lem:finite-flat-norm-base-change}
Let $w$ be admissible for $h$, $c$ as in Lemma~\ref{lem:nf-line-curve}(2), and $u$ any function on $Y^\an$; put
$u''=u\circ r\circ p$ on $(\Spec R')^\an$, where $r$ denotes restriction to $A$. If $z\in(\Spec Q)^\an$ satisfies
$|c\,\iota(h)|_z\ne0$, then $s':=(\iota^*z)|_B\in S_0^\an$ and $\Norm_qu''(z)=\Norm_\pi u(s')$.
\end{lemma}

\begin{proof}
$|h|_{s'}=|\iota(h)|_z\ne0$. Since $c$ is invertible in $\mathcal H(z)$, $\mathcal H(z)\otimes_QR'\cong\mathcal H(z)\otimes_QA_C$
compatibly with the maps from $A_F$ (Lemma~\ref{lem:nf-line-curve}(2)); so the points of $\Spec R'$ over $z$, with their
multiplicities and their restrictions to $A_F$, are those of $\Spec A_C$ over $z$. Now $A_C=Q\otimes_{B_F}A_F=Q\otimes_BA$
is the base change of $B\to A$ along $B\to Q$, and for every function the weighted fibre sum of a base change at $z$
equals the weighted fibre sum of $B\to A$ at the restriction $s'$ of $z$ to $B$
(Lemmas~\ref{spectrum:finite-base-change-multiplicity-mass} and~\ref{spectrum:all-function-finite-base-change}). Over
$S_0$, the multiplicities of Definition~\ref{def:nf-weights} are the fibre multiplicities of $B\to A$.
\end{proof}

\subsection{The weighted maximum principle on pieces}\label{subsec:nf-pieces}

In this subsection $F$ is a complete, algebraically closed, nontrivially valued field, $Q=F[t]_p$ with $p(0)\ne0$, $\xi$
is the point $t=0$ and $\zeta_\epsilon$ the Gauss points of Lemma~\ref{lem:nf-gauss-ray}. Let $R'=\prod_{i\in I}D_i$ be
a finite product of domains $D_i$, each containing $Q$ and finite over it, such that $R'$ is a free $Q$-module of rank
$d$. Let $x_1,\dots,x_m$ be the points of $(\Spec R')^\an$ over $\xi$ (they are $F$-points), with multiplicities $e_l$,
so $\sum_le_l=d$, and $q$ the map to $(\Spec Q)^\an$.

\begin{lemma}[Fibres over small discs]\label{lem:nf-fibre-shrink}
\lean{EnvelopeProperty.Berk.eventually_preimage_subset_of_continuous, EnvelopeProperty.eventually_nonflatFibre_nonflatGaussPoint_subset}
\leanok
\uses{spectrum:integral-compact-preimage, spectrum:topology, lem:nf-gauss-ray}
If $V$ is an open neighbourhood of $q^{-1}(\xi)$, then $q^{-1}(\zeta_\epsilon)\subseteq V$ for all sufficiently small
$\epsilon>0$.
\end{lemma}

\begin{proof}
Otherwise there are $\epsilon_n\to0$ and $x'_n\in q^{-1}(\zeta_{\epsilon_n})\setminus V$. The set
$K=\{\zeta_\epsilon:0\le\epsilon\le1\}$ is compact, as the continuous image of $[0,1]$, so $q^{-1}(K)$ is compact
(Lemma~\ref{spectrum:integral-compact-preimage}), and the $x'_n$ have a cluster point $x^*$ in the closed set
$q^{-1}(K)\setminus V$. By continuity $q(x^*)$ is a cluster point of $\zeta_{\epsilon_n}\to\zeta_0=\xi$, so
$q(x^*)=\xi$, the spaces being Hausdorff; thus $x^*\in q^{-1}(\xi)\subseteq V$, a contradiction.
\end{proof}

\begin{lemma}[Pieces of fibres over Gauss points]\label{lem:nf-pieces}
\lean{EnvelopeProperty.exists_nonflatPieces, EnvelopeProperty.nonflat_piece_norm_estimate, EnvelopeProperty.exists_nonflat_separating, EnvelopeProperty.exists_val_sub_eq_zero_of_mem_fibre_nonflatEvalPoint, EnvelopeProperty.Berk.val_algebraNorm_eq_prod_flatFibreMult}
\leanok
\uses{lem:nf-fibre-shrink, lem:nf-gauss-ray, spectrum:finite-free-scalar-norm}
There are pairwise disjoint open neighbourhoods $U_l$ of $x_l$ and $\epsilon_1>0$ with the following properties, where
$G_l(\epsilon)=q^{-1}(\zeta_\epsilon)\cap U_l$ for $0<\epsilon<\epsilon_1$:
\begin{enumerate}
\item $q^{-1}(\zeta_\epsilon)$ is the disjoint union of the $G_l(\epsilon)$;
\item $\sum_{x'\in G_l(\epsilon)}m_{x'}=e_l$;
\item $\prod_{x'\in G_l(\epsilon)}|b|_{x'}^{m_{x'}}\ge|b(x_l)|^{e_l}$ for every $b\in R'$;
\item for open neighbourhoods $V_l\subseteq U_l$ of $x_l$, $G_l(\epsilon)\subseteq V_l$ for all sufficiently small
$\epsilon$.
\end{enumerate}
\end{lemma}

\begin{proof}
Let $i(l)$ be the index with $x_l$ a point of $\Spec D_{i(l)}$, and $\varepsilon_i\in R'$ the idempotent of the factor
$D_i$. By the Chinese remainder theorem choose $\tilde e_l\in R'$ with $\tilde e_l(x_l)=1$ and $\tilde e_l(x_{l'})=0$ for
$l'\ne l$, and put $e'_l=\varepsilon_{i(l)}\tilde e_l$; it has the same values at the $x_{l'}$. Let
\[
U_l=\{|1-e'_l|<1\}\cap\bigcap_{l'\ne l}\{|e'_{l'}|<1\}.
\]
These sets are open and contain $x_l$; they are pairwise disjoint, since $|e'_l|<1$ and $|1-e'_l|<1$ cannot hold at
the same point. A point $x'\in U_l$ has $|e'_l|_{x'}=1$, hence $|\varepsilon_{i(l)}|_{x'}\ne0$ and, as
$\varepsilon_{i(l)}(1-\varepsilon_{i(l)})=0$, $|1-\varepsilon_{i(l)}|_{x'}=0$: so $x'$ is a point of $\Spec D_{i(l)}$. By
Lemmas~\ref{lem:nf-fibre-shrink} and~\ref{lem:nf-gauss-ray} choose $\epsilon_1\le\epsilon_0$ with
$q^{-1}(\zeta_\epsilon)\subseteq\bigcup_lU_l$ for $\epsilon<\epsilon_1$; this gives (1), and (4) follows in the same way
from Lemma~\ref{lem:nf-fibre-shrink} applied to $\bigcup_lV_l$.

Fix $\epsilon<\epsilon_1$ and $l$, write $G_{l'}=G_{l'}(\epsilon)$, $w_{l'}=\sum_{x'\in G_{l'}}m_{x'}$; then
$\sum_{l'}w_{l'}=d$, the module $R'$ being free of rank $d$. A point $x'\in G_l$ has zero kernel on $D_{i(l)}$: the
kernel is a prime of $D_{i(l)}$ lying over the zero kernel of $\zeta_\epsilon$, and $D_{i(l)}$ is an integral extension
of the domain $Q$. We claim: if $b\in R'$, $b(x_l)\ne0$, and $\kappa\in F$ satisfies $|\kappa|\ge1$ and
$|\kappa|>|b|_{x''}$ for all $x''\in q^{-1}(\zeta_\epsilon)$, then
\[
\prod_{x'\in G_l}|b|_{x'}^{m_{x'}}\cdot|\kappa|^{d-w_l}\ \ge\ |b(x_l)|^{e_l}\,|\kappa|^{d-e_l}. \tag{$\ast$}
\]
Since $b(x_l)\ne0$, $b$ is nonzero in $D_{i(l)}$, so $\beta=\min_{G_l}|b|>0$; also $\theta=\max_{G_l}|1-e'_l|<1$.
Choose $N\ge1$ with $|\kappa|\theta^N<\beta$, put $E=1-(1-e'_l)^N$ and $b_2=bE+\kappa(1-E)\in R'$. At $x'\in G_l$:
$|1-E|_{x'}\le\theta^N<1$, so $|E|_{x'}=1$ and $|bE|_{x'}=|b|_{x'}\ge\beta>|\kappa(1-E)|_{x'}$; hence
$|b_2|_{x'}=|b|_{x'}$. At $x'\in G_{l'}$, $l'\ne l$: $|e'_l|_{x'}<1$, and $E=\sum_{j=1}^N(-1)^{j+1}\binom Nje_l'^{\,j}$, so
$|E|_{x'}<1$, $|1-E|_{x'}=1$ and $|\kappa(1-E)|_{x'}=|\kappa|>|b|_{x'}\ge|bE|_{x'}$; hence $|b_2|_{x'}=|\kappa|$. Further
$E(x_l)=1$ and $E(x_{l'})=0$, so $b_2(x_l)=b(x_l)$ and $b_2(x_{l'})=\kappa$. The norm $\Norm_{R'/Q}(b_2)\in Q$ satisfies
$|\Norm_{R'/Q}(b_2)|_y=\prod_{x\mapsto y}|b_2|_x^{m_x}$ at every point $y$ of $(\Spec Q)^\an$
(Lemma~\ref{spectrum:finite-free-scalar-norm}), and $|\Norm(b_2)(0)|\le|\Norm(b_2)|_{\zeta_\epsilon}$
(Lemma~\ref{lem:nf-gauss-ray}). Thus
$|b(x_l)|^{e_l}|\kappa|^{d-e_l}\le\prod_{G_l}|b|^{m}\cdot|\kappa|^{\sum_{l'\ne l}w_{l'}}$, which is $(\ast)$.

(2) As $F$ is nontrivially valued, there is $\kappa\in F$ with $|\kappa|>1$. Apply $(\ast)$ to $b=1$: $|\kappa|^{d-w_l}\ge|\kappa|^{d-e_l}$, so
$w_l\le e_l$. Since $\sum_lw_l=d=\sum_le_l$, $w_l=e_l$ for all $l$. (3) If $b(x_l)=0$ there is nothing to prove;
otherwise apply $(\ast)$ with any admissible $\kappa$ and use $w_l=e_l$.
\end{proof}

\begin{lemma}[Weighted maximum principle for Fubini--Study pullbacks]\label{lem:nf-pieces-fs}
\lean{EnvelopeProperty.IsNonflatAwayFS, EnvelopeProperty.eventually_nonflat_weighted_maxPrinciple, EnvelopeProperty.eventually_nonflat_weighted_maxPrinciple_iInf, EnvelopeProperty.eventually_flatNormFun_nonflatEvalPoint_le}
\leanok
\uses{lem:nf-pieces}
Keep the notation of Lemma~\ref{lem:nf-pieces}, fix $l$ and $g\in R'$ with $g(x_l)\ne0$. Call a function $G$ on
$\{|g|\ne0\}\subseteq(\Spec R')^\an$ \emph{of FS type in $R'_g$} if $G=\max\{c_0,\max_i(a_i\log|\beta_i|+c_i)\}$ for
finitely many $\beta_i\in R'_g$, $a_i\in\Q_{>0}$ and $c_i\in\R$. There is $\epsilon_2>0$ such that for $0<\epsilon<\epsilon_2$:
\begin{enumerate}
\item $\sum_{x'\in G_l(\epsilon)}m_{x'}G(x')\ge e_l\,G(x_l)$ for every $G$ of FS type in $R'_g$;
\item if $W$ is the pointwise limit on $\{|g|\ne0\}$ of a decreasing sequence $G_j$ of functions of FS type in $R'_g$, then
$\sum_{x'\in G_l(\epsilon)}m_{x'}W(x')\ge e_l\,W(x_l)$.
\end{enumerate}
\end{lemma}

\begin{proof}
On the open neighbourhood $V_l=U_l\cap\{|g-g(x_l)|<|g(x_l)|\}$ of $x_l$ we have $|g|=|g(x_l)|$. By
Lemma~\ref{lem:nf-pieces}(4) choose $\epsilon_2$ with $G_l(\epsilon)\subseteq V_l$ for $\epsilon<\epsilon_2$. For
$\beta=b/g^n$ with $b\in R'$, Lemma~\ref{lem:nf-pieces}(2),(3) give
$\prod_{G_l(\epsilon)}|\beta|^{m}=\prod_{G_l(\epsilon)}|b|^m/|g(x_l)|^{ne_l}\ge|b(x_l)|^{e_l}/|g(x_l)|^{ne_l}=|\beta(x_l)|^{e_l}$.
(1) For each $i$, using $G\ge a_i\log|\beta_i|+c_i$ and $\sum_{G_l(\epsilon)}m=e_l$,
\[
\sum_{x'\in G_l(\epsilon)}m_{x'}G(x')\ge a_i\log\prod_{G_l(\epsilon)}|\beta_i|^{m}+e_lc_i\ge e_l\bigl(a_i\log|\beta_i(x_l)|+c_i\bigr),
\]
and likewise $\sum m_{x'}G(x')\ge e_lc_0$; take the maximum over the terms (with $\log0=-\infty$). (2) For a decreasing
sequence, the finite weighted sums with nonnegative weights converge to the weighted sum of the limits, so (2) follows
from (1) by letting $j\to\infty$ on both sides.
\end{proof}

\subsection{Local approximation on the normalized curve}\label{subsec:nf-local}

\begin{lemma}[Local approximants at rigid points]\label{lem:nf-local-approx}
\lean{EnvelopeProperty.exists_nonflatLocalApprox, EnvelopeProperty.exists_nonflatLocalApprox_away, EnvelopeProperty.exists_nonflatLocalApprox_chart, EnvelopeProperty.exists_nonflatLocalApprox_of_lift, EnvelopeProperty.nonflat_exists_chart_denominator, EnvelopeProperty.nonflat_apply_ne_zero_of_mul}
\leanok
\uses{def:nf-chart-regularizable, def:nf-line, lem:nf-line-curve, lem:nf-charpoly, def:fs}
Let $u$ be regularizable on the charts of a blowup, with witness $s$. Fix $0\ne\sigma_0\in s$, put $e_u=\chi_{\sigma_0}(0)$
(Lemma~\ref{lem:nf-charpoly} for $B\subseteq A$), let $(F,\xi)$ be a rigid datum for some $s\in S^\an$ and let $w$ be
admissible for $h\,e_u$. Then for every point $x$ of $(\Spec R')^\an$ there is $g\in R'$ with $g(x)\ne0$
such that $u''=u\circ r\circ p$ is, on $\{|g|\ne0\}$, the limit of a decreasing sequence of functions of FS type in
$R'_g$ (Lemma~\ref{lem:nf-pieces-fs}).
\end{lemma}

\begin{proof}
Put $y=r(p(x))\in Y^\an$. \emph{Case $y\notin Z^\an$.} Choose $g_0$ with $|g_0|_y\ne0$ and
$u|_{Y_{g_0}^\an}=\lim_jG_j$, $G_j\in\FS(Y_{g_0})$ decreasing. Let $g\in R'$ be the image of $g_0$; then
$|g|_x=|g_0|_y\ne0$, and $A\to R'$ extends to $A_{g_0}\to R'_g$. For $x'$ with $|g|_{x'}\ne0$, the point $r(p(x'))$ lies
in $Y_{g_0}^\an$, and $G_j(r(p(x')))$ is obtained by replacing each $\log|f|$, $f\in A_{g_0}$, in the formula for $G_j$ by
$\log|f'|(x')$ with $f'$ the image of $f$ in $R'_g$. These functions are of FS type in $R'_g$ and decrease to $u''$.
\emph{Case $y\in Z^\an$.} Let $D=D_P$ be the factor of $R'$ containing $x$, $\mathfrak m_x\subseteq D$ the maximal ideal
of $x$, and $\mathcal O=D_{\mathfrak m_x}$, a discrete valuation ring with valuation $\ord_x$
(Lemma~\ref{lem:nf-line-curve}(1)). The image of $\sigma_0$ in $D$ is nonzero: by Lemma~\ref{lem:nf-charpoly}(3)
$e_u\in\sigma_0A$, so otherwise the image of $\iota(e_u)$ in $D$ would vanish, while $Q\to D$ is injective and
$\iota(e_u)\ne0$. Choose $\sigma\in s$ with nonzero image in $D$ and $\ord_x(\sigma)$ minimal among such elements. Then
$\sigma'/\sigma\in\mathcal O$ for every $\sigma'\in s$ (if $\sigma'$ maps to $0$ this is $0$). The chart $A[s/\sigma]$ is
contained in $A_\sigma$, and $A\to D$ extends to $A_\sigma\to D_\sigma\subseteq\mathrm{Frac}\,D$; the image of
$A[s/\sigma]$ is generated by the images of $A$ and of the $\sigma'/\sigma$, so it lies in $\mathcal O$, and as
$A[s/\sigma]$ is finitely generated it lies in $D_{g'}$ for some $g'\in D\setminus\mathfrak m_x$. Let $g\in R'$ be $g'$
on the factor $D$ and $0$ on the other factors; then $R'_g=D_{g'}$, $g(x)\ne0$, and we obtain
$\psi\colon A[s/\sigma]\to R'_g$ extending $A\to R'$. For $x'$ with $|g|_{x'}\ne0$, the restriction $\psi^*x'$ of $x'$ to
$A[s/\sigma]$ satisfies $\mu_\sigma(\psi^*x')=r(p(x'))$, so $u''(x')=(u\circ\mu_\sigma)(\psi^*x')$. Since
$u\circ\mu_\sigma\in\Rcal(A[s/\sigma])$, its Fubini--Study approximants composed with $\psi^*$ are as in the first case.
The second case uses $y\in Z^\an$ only through the nonvanishing of $\sigma_0$ on $D$, which holds on every factor; so it applies at
every point $x$ (if the kernel of $x$ on $D$ is zero, $\mathcal O=\mathrm{Frac}\,D$), and the covering hypothesis of the first case is not
needed. The first case is kept as a separate statement.
\end{proof}

\subsection{Identification and the nonflat norm}\label{subsec:nf-main}

\begin{proposition}[Identification off the \'etale locus]\label{prop:nf-identification}
\lean{EnvelopeProperty.IsChartRegularizable.eq_nonflatNorm, EnvelopeProperty.IsChartRegularizable.eq_nonflatNorm_of_satisfiesHullIneq, EnvelopeProperty.NonflatRigidEmbedding.eq_sum_rigidWeight_of_line, EnvelopeProperty.nonflat_eq_flatNormFun_of_gaussRay, EnvelopeProperty.NonflatRigidEmbedding.exists_nonflatLocalApprox_line}
\leanok
\uses{lem:nf-chart-regularizable-closure, def:nf-weights, lem:nf-line-exists, lem:nf-line-curve, lem:nf-line-weights, lem:nf-hull-approach, lem:nf-generic-fibre, lem:nf-pieces, lem:nf-pieces-fs, lem:nf-local-approx, lem:nf-general-gauss-usc-upper, lem:nf-general-order-squeeze, lem:nf-whole-rigid-weighted-norm}
Let $u$ be regularizable on the charts of a blowup. Let $v\in\Rcal(S)$ (it suffices that $v$ is upper semicontinuous and satisfies the hull inequality) with
$v=\Norm_\pi u$ on $S_0^\an$. Then
$v(s)=\Norm_\pi u(s)$ for every $s\in S^\an$.
\end{proposition}

\begin{proof}
Let $s\notin S_0^\an$, with $(F,\xi)=(F_s,\xi_s)$ and $b_s,\chi_s$ as in Definition~\ref{def:nf-weights}; let $e_u$ be as
in Lemma~\ref{lem:nf-local-approx} and $w$ a direction admissible for $h\,e_u$ (Lemma~\ref{lem:nf-line-exists}); it is
admissible for $h$. Let $Q,R',q,p$ be as in Definition~\ref{def:nf-line}, $c$ as in Lemma~\ref{lem:nf-line-curve}(2),
$u''=u\circ r\circ p$, and $x_1,\dots,x_m$ the points over $\xi$, with multiplicities $e_l$. The hypotheses of the
subsection on pieces hold by Lemma~\ref{lem:nf-line-curve}. By Lemma~\ref{lem:nf-line-weights},
\[
\Norm_qu''(\xi)=\sum_le_l\,u(r(p(x_l)))=\sum_zM_z\,u(r(z))=\sum_{\pi(y)=s}m_y(\pi)u(y)=\Norm_\pi u(s).
\]
\emph{Lower bound.} For each $l$ take $g_l$ from Lemma~\ref{lem:nf-local-approx} and $\epsilon_2^{(l)}$ from
Lemma~\ref{lem:nf-pieces-fs}(2). For $0<\epsilon<\min_l\epsilon_2^{(l)}$, Lemma~\ref{lem:nf-pieces}(1) and
Lemma~\ref{lem:nf-pieces-fs}(2) give
$\Norm_qu''(\zeta_\epsilon)=\sum_l\sum_{G_l(\epsilon)}m\,u''\ge\sum_le_lu''(x_l)=\Norm_qu''(\xi)$.
\emph{Upper bound.} $u''$ is upper semicontinuous, $u$ being so (Lemma~\ref{lem:nf-chart-regularizable-closure}(1)), and $u''(x_l)<+\infty$ (the local
approximants of Lemma~\ref{lem:nf-local-approx} are of FS type). For
reals $c_l>u''(x_l)$ the sets $V_l=U_l\cap\{u''<c_l\}$ are open neighbourhoods of $x_l$, so by
Lemma~\ref{lem:nf-pieces}(2),(4) $\Norm_qu''(\zeta_\epsilon)\le\sum_le_lc_l$ for small $\epsilon$. Hence
$\limsup_{\epsilon\to0}\Norm_qu''(\zeta_\epsilon)\le\Norm_qu''(\xi)$.
\emph{Comparison with $v$.} The point $\zeta_\epsilon$, $\epsilon>0$, has zero kernel, so $|c\,\iota(h)|_{\zeta_\epsilon}\ne0$
and Lemma~\ref{lem:nf-generic-fibre} gives $\Norm_qu''(\zeta_\epsilon)=\Norm_\pi u(s_\epsilon)=v(s_\epsilon)$, with
$s_\epsilon\in S_0^\an$ as in Lemma~\ref{lem:nf-hull-approach}. By that lemma, $v(s)\le v(s_\epsilon)$ for small $\epsilon$,
so $v(s)\le\limsup_\epsilon\Norm_qu''(\zeta_\epsilon)\le\Norm_\pi u(s)$; and
$v(s)\ge\limsup_\epsilon v(s_\epsilon)\ge\Norm_\pi u(s)$ by the lower bound. For $s\in S_0^\an$ the identity is the
hypothesis.
\end{proof}

\begin{proposition}[Nonflat norms for functions regularizable on blowup charts]\label{prop:nf-nonflat-norm}
\lean{EnvelopeProperty.IsChartRegularizable.isRegularizable_nonflatNorm, EnvelopeProperty.nonempty_normWeights_isChartRegularizable, EnvelopeProperty.ChartNonflatNorms.of_wuRegularization, EnvelopeProperty.hasLogGrowth_nonflatNorm, EnvelopeProperty.IsChartRegularizable.isRegularizable_awayComap_goodLocus}
\leanok
\uses{conv:nf-setting, def:nf-chart-regularizable, lem:nf-chart-regularizable-closure, lem:nf-good-locus, def:nf-weights, lem:nf-weights-basic, prop:gluing, lem:reg-closure-elementary, lem:finite-etale-norm, lem:growth, lem:nonflat-extension, prop:nf-identification}
\source{fhx-envelope-20261028:prop:nonflat-norm}
In the setting of Convention~\ref{conv:nf-setting}, the integers $m_y(\pi)$ of Definition~\ref{def:nf-weights} satisfy
$m_y(\pi)\ge1$ and $\sum_{\pi(y)=s}m_y(\pi)=d$ for every $s\in S^\an$. If $u$ is regularizable on the charts of a blowup,
of logarithmic growth, and every $x\in Y^\an\setminus Z^\an$ lies in some $Y_{g_0}^\an$ with $u|_{Y_{g_0}^\an}\in\Rcal(Y_{g_0})$,
then $\Norm_\pi u\in\Rcal(S)$.
\end{proposition}

\begin{proof}
The first assertion is Lemma~\ref{lem:nf-weights-basic}. The variety $Y_0$ is normal and affine and contained in
$Y\setminus Z$, and it is covered by the principal opens $Y_{\varphi(h)g_0}$ on which $u$ is regularizable
(Lemma~\ref{lem:reg-closure-elementary}(4)); $u$ is upper semicontinuous and of logarithmic growth, so
$u|_{Y_0^\an}\in\Rcal(Y_0)$ by Proposition~\ref{prop:gluing}(2). By Lemma~\ref{lem:finite-etale-norm} applied to the finite
free \'etale $B_h\to A_{\varphi(h)}$, $v_0=\Norm_{\pi|_{Y_0}}u\in\Rcal(S_0)$; for $s\in S_0^\an$ the fibre algebras of
$B\to A$ and $B_h\to A_{\varphi(h)}$ at $s$ coincide, so $v_0=\Norm_\pi u$ on $S_0^\an$. By logarithmic growth of $u$ and
Lemma~\ref{lem:growth}(2), $u\le C\,g\circ\pi+C'$ with $g\in\FS(S)$, $C\ge0$; as the weights over each point sum to $d$,
$v_0\le dC\,g+d|C'|$ on $S_0^\an$, and $dC\,g$ may be replaced by a function in $\FS(S)$ (enlarging $C$ to a rational
number and using $g\ge b_0$). Lemma~\ref{lem:nonflat-extension} gives $v\in\Rcal(S)$ with $v=v_0$ on $S_0^\an$, and
Proposition~\ref{prop:nf-identification} gives $v=\Norm_\pi u$.
\end{proof}

\begin{proposition}[Projection step for functions regularizable on blowup charts]\label{prop:nf-projection-step}
\lean{EnvelopeProperty.IsRegularizable.of_projection_chartRegularizable, EnvelopeProperty.ChartNonflatNorms, EnvelopeProperty.ChartNonflatNorms.of_nonflatNorms, EnvelopeProperty.NormWeights.ofSubclass}
\leanok
\uses{prop:projection-step, prop:nf-nonflat-norm, lem:nf-chart-regularizable-closure}
\source{fhx-envelope-20261028:prop:projection-step}
Proposition~\ref{prop:projection-step} holds for every $u$ that is regularizable on the charts of a blowup (in place of
weakly psh), with the weights of Proposition~\ref{prop:nf-nonflat-norm}.
\end{proposition}

\begin{proof}
The proof of Proposition~\ref{prop:projection-step} uses of the weakly psh hypothesis only upper semicontinuity, stability
of the class under adding and taking maxima with regularizable functions, stability under $u\mapsto u\circ p+\log|w|$ on
$\widetilde Y$, and Proposition~\ref{prop:nonflat-norm} for $\pi$ and for $\pi\times\mathrm{id}$ on $\widetilde Y$, whose
base $S\times\A^1$ is a principal open of $\A^{n+1}$ and whose map is again finite, injective and generically \'etale of
degree $d$. The first three are Lemma~\ref{lem:nf-chart-regularizable-closure}(1),(2),(4), and the last is
Proposition~\ref{prop:nf-nonflat-norm} applied to $\pi$ and to $\pi\times\mathrm{id}$.
\end{proof}

\begin{theorem}[Local regularization for functions regularizable on blowup charts]\label{thm:nf-local-regularization}
\lean{EnvelopeProperty.IsRegularizable.of_isChartRegularizable, EnvelopeProperty.IsRegularizable.of_isChartRegularizable_of_isAlgClosed, EnvelopeProperty.IsRegularizable.of_isChartRegularizable_of_convexPLApproximation, EnvelopeProperty.IsRegularizable.of_isChartRegularizable_of_localPositivity}
\leanok
\uses{thm:normal-local-regularization, prop:nf-projection-step, lem:nf-chart-regularizable-closure}
\source{fhx-envelope-20261028:thm:normal-local-regularization}
Theorem~\ref{thm:normal-local-regularization} holds for every $u$ that is regularizable on the charts of a blowup, of
logarithmic growth, and regularizable on principal opens covering $Y\setminus Z$.
\end{theorem}

\begin{proof}
The proof of Theorem~\ref{thm:normal-local-regularization} restricts $u$ to principal opens $Y_h$ and applies the
projection step there. Restrictions stay in the class (Lemma~\ref{lem:nf-chart-regularizable-closure}(3)), and
Proposition~\ref{prop:nf-projection-step} replaces Proposition~\ref{prop:projection-step}.
\end{proof}

\begin{corollary}[Psh descent without the nonflat norm hypothesis]\label{cor:nf-psh-descent}
\lean{EnvelopeProperty.ProjSubvariety.isPSH_of_isPSH_comp_proj, EnvelopeProperty.ProjSubvariety.isPSH_of_isPSH_comp_proj_of_convexPLApproximation}
\leanok
\uses{thm:pg-psh-descent, lem:nf-descent-chart, thm:nf-local-regularization}
\source{fhx-envelope-20261028:cor:normal-weak-global}
Theorem~\ref{thm:pg-psh-descent} holds without assuming Proposition~\ref{prop:nonflat-norm}.
\end{corollary}

\begin{proof}
In the proof of Theorem~\ref{thm:pg-psh-descent}, Theorem~\ref{thm:normal-local-regularization} is applied to the chart
functions $u_i$, which are regularizable on the charts of a blowup by Lemma~\ref{lem:nf-descent-chart}; apply
Theorem~\ref{thm:nf-local-regularization} instead.
\end{proof}

\begin{remark}[Relation to the source]\label{rmk:nf-source}
\uses{prop:nonflat-norm, lem:curves, prop:nf-nonflat-norm}
The proof of \cite[Proposition~5.22]{FHX} chooses one line through the rigid point for each $s$, defines the weights by
the multiplicities of the normalized curve over that line, and identifies the extended norm with $\Norm_\pi u$ through
\cite[Lemma~5.21]{FHX}: subharmonicity of norms along finite flat maps of curves, equality of subharmonic functions that
agree off finitely many rigid points, and subharmonicity of pullbacks of weakly psh functions to curves, all from
\cite[Proposition~3.5(1),(4),(6),(8) and Proposition~4.2]{GR25}. Here the weights are defined by characteristic
polynomials; by Lemma~\ref{lem:nf-line-weights} they coincide with the line multiplicities of \cite{FHX} for every line
admissible for $h$, so the line may be chosen after the function $u$ (it must avoid the image of the blowup centre,
through the element $e_u$). The identification uses the hull inequality along Gauss points of discs in the line
(Lemma~\ref{lem:nf-hull-approach}) instead of the maximum principle for subharmonic functions, and the weighted inequality
of Lemma~\ref{lem:nf-pieces-fs} on the normalized curve, where the function is approximated by Fubini--Study functions of
a blowup chart (Lemma~\ref{lem:nf-local-approx}). Only the discrete valuation rings of the normalized curve are used, not
its smoothness. The rigid field is taken nontrivially valued (Definition~\ref{def:nf-rigid-data}); the completed
algebraic closure of $\mathcal H(s)$ used in \cite{FHX} is trivially valued when $\mathcal H(s)$ is. The class of
Definition~\ref{def:nf-chart-regularizable} is smaller than the class of weakly psh functions of \cite[Proposition~5.22]{FHX}
and contains the functions of \cite[Corollary~5.26]{FHX} (Lemma~\ref{lem:nf-descent-chart}); the statement of
Proposition~\ref{prop:nonflat-norm} for all weakly psh functions is not proved by this route, and Lemma~\ref{lem:curves}
is not used.
\end{remark}
